\clearpage
\chapter*{GLOSARIO}
\addcontentsline{toc}{chapter}{GLOSARIO}

\begin{description}
  
  \item[API] \textit{Application Programming Interface.} Conjunto de definiciones y protocolos que permiten la comunicación entre diferentes aplicaciones o servicios.
  
  \item[BLE] \textit{Bluetooth Low Energy.} Versión del protocolo Bluetooth diseñada para aplicaciones de bajo consumo energético, ideal para dispositivos IoT.
  
  \item[CPS] \textit{Cyber-Physical Systems.} Sistemas que integran componentes computacionales con procesos físicos, permitiendo la interacción y control en tiempo real.
  
  \item[CSV] \textit{Comma-Separated Values.} Formato de archivo de datos tabulares donde los valores están separados por comas, utilizado para el almacenamiento de registros.
  
  \item[EMI] \textit{Electromagnetic Interference.} Interferencia electromagnética que afecta el funcionamiento de dispositivos electrónicos, generada por fuentes externas o internas.
  
  \item[EMO] \textit{Efectividad Mecánica Operativa.} Métrica adaptada del OEE que mide la eficiencia mecánica del equipo, calculada como el producto de la disponibilidad y el rendimiento.
  
  \item[ESP32] Microcontrolador de bajo costo con conectividad Wi-Fi y Bluetooth integrada, desarrollado por Espressif Systems, ampliamente utilizado en aplicaciones IoT.
  
  \item[FSM] \textit{Finite State Machine.} Modelo matemático que representa sistemas con un número finito de estados y transiciones entre ellos, utilizado para el control secuencial de procesos.
  
  \item[GPIO] \textit{General Purpose Input/Output.} Pines de propósito general en un microcontrolador que pueden configurarse como entrada o salida digital.
  
  \item[HMI] \textit{Human-Machine Interface.} Interfaz que permite la interacción entre el usuario y un sistema industrial, facilitando el monitoreo y control.
  
  \item[I2C] \textit{Inter-Integrated Circuit.} Protocolo de comunicación serie síncrono para conectar dispositivos periféricos a un microcontrolador mediante un bus de dos hilos.
  
  \item[IIoT] \textit{Industrial Internet of Things.} Aplicación del Internet de las Cosas en entornos industriales para la monitorización y control de procesos.
  
  \item[IoT] \textit{Internet of Things.} Red de dispositivos físicos interconectados que recopilan y comparten datos a través de Internet.
  
  \item[IP54] Grado de protección internacional que indica resistencia al polvo y a salpicaduras de agua en equipos electrónicos.
  
  \item[JSON] \textit{JavaScript Object Notation.} Formato de texto ligero para el intercambio de datos, basado en la sintaxis de objetos de JavaScript.
  
  \item[MQTT] \textit{Message Queuing Telemetry Transport.} Protocolo de mensajería ligero de publicación/suscripción, ampliamente utilizado en entornos IoT e IIoT.
  
  \item[MTTR] \textit{Mean Time To Repair.} Tiempo medio necesario para reparar un sistema o componente después de un fallo, medido desde la detección hasta la restauración.
  
  \item[Node-RED] Herramienta de programación visual basada en flujos para el desarrollo de aplicaciones IoT, SCADA y automatización industrial.
  
  \item[OASIS] \textit{Organization for the Advancement of Structured Information Standards.} Consorcio internacional que desarrolla estándares para la interoperabilidad de sistemas.
  
  \item[OEE] \textit{Overall Equipment Effectiveness.} Indicador clave de rendimiento que mide la eficiencia global de un equipo, calculado como el producto de disponibilidad, rendimiento y calidad.
  
  \item[PC817] Optoacoplador utilizado para aislar galvánicamente circuitos eléctricos, protegiendo las entradas digitales de microcontroladores.
  
  \item[PLC] \textit{Programmable Logic Controller.} Controlador lógico programable utilizado en la automatización industrial para el control de procesos y maquinaria.
  
  \item[PYMES] Pequeñas y Medianas Empresas, sector productivo que constituye la mayor parte del tejido empresarial latinoamericano.
  
  \item[QoS] \textit{Quality of Service.} Calidad de Servicio, mecanismo que garantiza la entrega de mensajes en protocolos de comunicación como MQTT.
  
  \item[SCADA] \textit{Supervisory Control and Data Acquisition.} Sistema de supervisión, control y adquisición de datos para la monitorización de procesos industriales.
  
  \item[SQLite] Sistema de gestión de bases de datos relacional ligero, utilizado para el almacenamiento local de datos en aplicaciones embebidas.
  
  \item[TDD] \textit{Test-Driven Development.} Metodología de desarrollo de software donde las pruebas se escriben antes del código, guiando la implementación.
  
  \item[TIC] Tecnologías de la Información y Comunicación, conjunto de tecnologías que permiten el procesamiento, almacenamiento y transmisión de información.
  
  \item[Webots] Plataforma de simulación 3D para robótica y automatización, que permite modelar entornos físicos y programar controladores en Python.
  
  \item[WiFi] Tecnología de comunicación inalámbrica que permite la conexión de dispositivos a redes de área local a través de ondas de radio.
  
  \item[\textit{Legacy}] Término utilizado para describir sistemas, equipos o software que son obsoletos pero que aún están en operación.
  
  \item[\textit{Retrofit}] Proceso de modernización de equipos o sistemas existentes mediante la incorporación de nuevas tecnologías o componentes.
  
  \item[\textit{Bypass} Ciberfísico] Arquitectura no intrusiva que permite la integración de sistemas modernos con máquinas \textit{legacy} sin modificar su cableado o lógica original.
  
  \item[\textit{Snubber}] Circuito eléctrico utilizado para suprimir picos de tensión y transitorios en sistemas inductivos, protegiendo componentes electrónicos.
  
  \item[\textit{Watchdog}] Mecanismo de supervisión en sistemas embebidos que reinicia automáticamente el dispositivo en caso de bloqueo o fallo del firmware.
  
  \item[\textit{Virtual Commissioning}] Proceso de verificación y validación de sistemas de control en entornos virtuales antes de su implementación en hardware real.
  
  \item[\textit{Soft real-time}] Sistema de control donde los plazos de tiempo no son estrictos pero se busca minimizar la latencia para mantener la sincronización.
  
  \item[\textit{Brownfield}] Término utilizado para describir instalaciones industriales o sistemas existentes que serán modernizados o ampliados.
  
  \item[\textit{Edge Computing}] Modelo de procesamiento de datos que se ejecuta cerca de la fuente de generación de datos, reduciendo la latencia y el ancho de banda.
  
  \item[\textit{Cloud Computing}] Modelo de computación que permite el acceso a recursos informáticos compartidos a través de Internet, bajo demanda.
  
\end{description}

